\chapter{Методология исследования}

\section{Исследовательская позиция и preregistration}
Исследователь фиксирует предположения, процедуры, критерии допуска и границы
вывода отдельно от результатов. Изменения после соответствующей точки
заморозки не переписывают исходный подтверждающий анализ.

\section{Статистическая единица и независимость}
Для Stage~1/2 и соответствующих Stage~3 сравнений независимой статистической
единицей является независимо обученная модель, идентифицируемая
\texttt{model\_seed}. Пакеты, слои, состояния и отдельные измерения одной
модели не трактуются как независимые репликации.

\section{Разделение validation, test и post-action oracle}
Pilot использует validation. Доступ к test разрешается только после
соответствующей заморозки. В QWake post-action oracle используется для аудита
решений после формирования pre-action представления и не входит в допустимый
вход policy.

\section{Последовательность экспериментальных стадий}
Методология разделяет Stage~1/2, Stage~3A, Stage~3B mechanism-attribution и
QWake C1/C2. Для каждой стадии фиксируются собственные входы, критерии
корректности, вычислительная среда и разрешённые выводы.

\section{Эквивалентность, различие и отрицательный результат}
Отсутствие обнаруженного различия не считается эквивалентностью без отдельного
equivalence анализа. Аналогично, bounded-negative QWake C2 не расширяется до
универсального утверждения о невозможности безопасного или экономически
полезного recognizer вне замороженного candidate family и cost accounting.

\section{Provenance и воспроизводимость}
Git commit, immutable environment/image identities, request/authorization
bindings, manifests, receipts и checksums используются для привязки
последовательных стадий. Диссертационный PDF собирается отдельным CI workflow
из Git-tracked source и thesis-facing evidence summaries; сборка текста не
изменяет научные evidence artifacts.
